\documentclass[conference]{IEEEtran}
\usepackage{amsmath,amssymb}
\usepackage{booktabs}
\usepackage{graphicx}
\usepackage{url}
\usepackage[hidelinks]{hyperref}
\usepackage{microtype}

\begin{document}

\title{Event-Gated Ensemble Enhancement with ASR-Consistent Training for Indic Noise Event Removal}

\author{\IEEEauthorblockN{Shivam Dwivedi, Animish Yadav, Shubhi Jain}
\IEEEauthorblockA{Team TDS, Indraprastha Institute of Information Technology Delhi (IIIT Delhi), India\\
\{shivam21352, animish23089, shubhi23517\}@iiitd.ac.in}}

\maketitle

\begin{abstract}
We present our Track~2 system for the Datathon@IndoML~2026 noise event removal challenge. The task is to
suppress non-speech events (horns, animals, breathing, children, alarms, machinery, music) in real-world Indic
speech from the Vaani corpus. The score is the SI-SDR on synthetic clips plus $100\times$ the WER reduction of a
frozen mandated ASR (SraVaani-1.0) over all clips. A single learned enhancer is applied identically to every clip.
It averages a waveform-domain DNS64 (Demucs) model and a band-split RNN with a learned frame-level
\emph{event gate}, so the network detects where events are and leaves noise-free speech untouched. Both models are
trained on on-the-fly mixtures whose event statistics follow the validation set, with an
\emph{ASR-consistency loss} on the frozen SraVaani encoder. A second stage fine-tunes them on longer, denser and
louder mixtures, selecting checkpoints on a \emph{pseudo-test} assembled from held-out validation clips. That
pseudo-test predicted our leaderboard score within 0.13. Our final submission (Codabench 995085) reaches
SI-SDR 12.78\,dB and $\Delta$WER $+1.87$ (synthetic $+5.41$), for a combined score of 14.65 (Phase~2 rank~13), with
the hidden ASR audit passed.
\end{abstract}

\section{Training Data and Preprocessing}
\textbf{Vaani Noise Event Dataset}~\cite{vaani} (90,637 segments, 154.6\,h, Gold/Silver/Bronze annotation tiers).
We use it as a clean-speech pool. From Gold and Silver clips we keep non-event regions of at least 1.5\,s that lie
at least 0.25\,s from any labelled event, and only from clips whose 10th-percentile frame energy is more than
28\,dB below the clip RMS. This yields 8\,h of low-background speech.

\textbf{IndoML validation set} (2,501 clips; 1,806 natural, 695 synthetic with clean references). The 695 synthetic
pairs are split 560 for training and 135 for a holdout. From the 560 we take the clean references (0.81\,h) and
\emph{1,783 noise events} cut from $\text{noisy}-\text{clean}$ at the labelled spans ($\pm$50\,ms). In these pairs
the added noise lies entirely inside the spans; outside them the noisy signal equals the clean one.

\textbf{External data.} We use ESC-50~\cite{esc50} (2,000 environmental clips, resampled to 16\,kHz and trimmed to
the region within 30\,dB of the peak) as a second noise source, so the model generalises beyond the validation
noise. We also use Google FLEURS~\cite{fleurs} Indic train shards (hi, bn, te, ta, kn, ml, mr, or, as, gu, pa, ne;
at most 600 utterances per language, 12\,s or shorter) as additional clean read speech. In total the banks hold
3,783 noise events (2.72\,h) and up to 17.6\,h of clean speech, stored as memory-mapped float16 arrays.

\textbf{Augmentation pipeline} (on the fly; every batch is new).
\begin{enumerate}
\item Clean speech is a validation reference with probability 0.5, otherwise a Vaani or FLEURS segment.
Utterances, each rescaled within $\pm$3\,dB and separated by 50--600\,ms gaps, are joined until the 6--10\,s crop
is full. The crop is scaled to $-30\ldots-15$\,dBFS.
\item The number of events is $k=n\cdot U(1,L/4.5\,\text{s})\cdot U(\text{ev\_mult})$, where $n$ is drawn from the
per-clip event counts of the validation synthetic set.
\item Each event is a validation noise event ($p{=}0.7$) or an ESC-50 crop of empirical event duration ($p{=}0.3$)
with 5\,ms fades and random polarity, placed at a uniform offset. Its level relative to the clean RMS is drawn
from the empirical validation distribution (median $+4.3$\,dB), plus $\mathcal{N}(0,2)$\,dB, plus
$U(\text{db\_shift})$.
\item The frame event mask (hop 128) is the target of the activity head. Mixture and target are peak-normalised
jointly.
\end{enumerate}
The simulated mixtures match validation (median input SI-SDR $-2.5$ vs.\ $-2.6$\,dB; event coverage 43\% vs.\ 45\%).
Stage~2 uses $\text{ev\_mult}=(1,2)$, $\text{db\_shift}=(0,3)$\,dB and 10\,s crops.

\section{Proposed Method}
\subsection{Models}
\textbf{DNS64.} The Demucs waveform U-Net with an LSTM bottleneck~\cite{demucs} (hidden width 64, 33.5\,M
parameters), initialised from the public checkpoint \texttt{dns64-a7761ff99a7d5bb6.th} (facebookresearch/denoiser).

\textbf{Event-gated BSRNN} (ours, 11.3\,M parameters), after band-split RNN~\cite{bsrnn}.
\begin{itemize}
\item A 512-point STFT (hop 128) is split into 37 bands of $16{\times}4$, $16{\times}8$, $4{\times}16$ and 1 bins.
Each band is normalised and projected to 112 features.
\item Six blocks alternate residual BLSTMs over time and over bands. Per-band GLU MLPs output a complex mask
$\mathbf{M}$, so $\hat{s}=\text{iSTFT}(\mathbf{M}\odot\mathbf{X})$.
\item A gate head maps the concatenated band features to an event probability $p(t)$. The output is
\begin{equation}
y = x + g(t)\,(\hat{s}-x),\qquad g(t)=\gamma+(1-\gamma)\,p(t),
\end{equation}
with floor $\gamma=0.2$, so every frame is still processed while event-free speech stays close to the input.
\end{itemize}

\subsection{Losses}
For mixture $x$, clean target $s$ and event mask $m$,
\begin{align}
\mathcal{L}_{\text{BSRNN}} ={}& \tfrac12\mathcal{L}_{\text{rec}}(\hat{s},s)+\tfrac12\mathcal{L}_{\text{rec}}(y,s)
+ \text{BCE}(p,m) + \lambda\,\mathcal{L}_{\text{ASR}},\\
\mathcal{L}_{\text{rec}}(a,s) ={}& -\text{SI-SDR}(a,s) + 0.5\,\text{MRSTFT}(a,s),
\end{align}
and $\mathcal{L}_{\text{DNS64}}=\mathcal{L}_{\text{rec}}(\hat{s},s)+\lambda\mathcal{L}_{\text{ASR}}$.
\begin{itemize}
\item MRSTFT is the L1 distance between log-magnitudes at FFT sizes 256, 512 and 1024~\cite{pwg}.
\item SI-SDR is the metric itself~\cite{sisdr}.
\item The \emph{ASR-consistency loss} $\mathcal{L}_{\text{ASR}}$ is the L1 distance between the frozen SraVaani-1.0
encoder~\cite{sravaani} outputs on the enhanced and clean signals, normalised by the target's mean magnitude. It
is applied every second step, with $\lambda=0.1$ (0.2 for one DNS64 run). Its front end replicates the ASR's
pre-emphasis, 128-bin log-mel and per-feature normalisation. Gradients flow through the frozen encoder into the
enhancer, so the enhancer learns outputs the mandated ASR transcribes like clean speech.
\end{itemize}

\subsection{Training strategy}
\textbf{Optimiser.} AdamW~\cite{adamw} (weight decay $10^{-2}$) with a 500-step warm-up and a cosine schedule driven by
\emph{wall-clock time}, which guarantees completion inside a fixed GPU allocation. We use gradient clipping at 5
and bf16 autocast for BSRNN, and keep the best checkpoint by evaluation SI-SDR.

\textbf{Stage 1} (holdout selection):
\begin{itemize}
\item DNS64: 6\,s crops, learning rate $10^{-4}$, 210\,min on an NVIDIA B200.
\item BSRNN: trained from scratch, 6\,s crops, learning rate $10^{-3}$, 180\,min on an A100 80\,GB.
\item A second DNS64: 8\,s crops, learning rate $2{\cdot}10^{-4}$, $\lambda=0.2$, 150\,min on an A100.
\end{itemize}

\textbf{Pseudo-test.} After the first submission showed a large gap between validation and test, we built a
pseudo-test only from the 135 holdout clips. It has 106 clips of 2--3 joined holdout references (median 13.5\,s)
carrying their own noise plus $0.5\times$ extra holdout events at $+1.5$\,dB. Its median input SI-SDR is
$-4.5$\,dB, which matches the noise level we estimated for the real synthetic test clips ($-4.48$\,dB). It reproduces
the real per-clip score distribution of our first submission (median 12.79 vs.\ 12.27\,dB; 5th percentile 4.78
vs.\ 4.59\,dB).

\textbf{Stage 2} fine-tunes DNS64 and BSRNN on test-like mixtures for 50\,min each on the B200. Checkpoints are
selected on the pseudo-test, and BSRNN's clean pool excludes the Vaani segments.

\textbf{Inference.} The final system averages the two stage-2 models, with gate floor 0.2. Output RMS is matched to
the input, clips are processed in 20\,s chunks with a 2\,s cross-fade, and audio is written as PCM-16. Transcripts
come from the unmodified SraVaani-1.0 (fp32, greedy TDT~\cite{tdt}) on the written WAVs.

\textbf{Candidate selection.} Candidates are ranked by an \emph{estimated leaderboard score}:
\begin{itemize}
\item pseudo-test SI-SDR minus a 0.70\,dB calibration offset;
\item plus $0.32\times$ the synthetic $\Delta$WER, measured on holdout clips against ASR-on-clean transcripts and
calibrated to our first submission;
\item plus $0.68\times$ the natural-validation $\Delta$WER.
\end{itemize}
The weights 0.32 and 0.68 are the synthetic and natural shares of words in the test set.

\section{Results and Conclusion}
\begin{table}[t]
\centering
\caption{SI-SDR (dB) on the validation holdout (135 clips) and the pseudo-test (106 clips), and the Phase~2 test
score (SI-SDR / $\Delta$WER / combined).}
\label{tab:res}
\footnotesize
\setlength{\tabcolsep}{3.5pt}
\begin{tabular}{lccc}
\toprule
System & Holdout & Pseudo-test & Test (Codabench)\\
\midrule
Noisy input & $-0.76$ & $-3.75$ & --\\
DNS64 (stage 1) & 19.04 & 12.81 & --\\
BSRNN, gated (stage 1) & 19.04 & 12.67 & --\\
DNS64 + BSRNN & 19.78 & 13.36 & 12.66 / 1.80 / 14.46\\
+ second DNS64 & 19.71 & 13.36 & -- / -- / 14.41\\
\textbf{Stage-2 DNS64 + BSRNN} & -- & \textbf{13.52} & \textbf{12.78 / 1.87 / 14.65}\\
\bottomrule
\end{tabular}
\end{table}

Table~\ref{tab:res} summarises the results.
\begin{itemize}
\item \textbf{Final submission (995085):} SI-SDR 12.78\,dB and $\Delta$WER $+1.87$, of which $+5.41$ on synthetic and
$+0.21$ on natural clips. Combined 14.65; the hidden ASR audit passed with 0\% mismatch.
\item \textbf{Architectural diversity helps.} DNS64 + BSRNN gains 0.75\,dB over either model alone and turns the
holdout $\Delta$WER positive, whereas averaging two DNS64 models adds nothing.
\item \textbf{The learned gate is only useful with good data.} BSRNN gated at $\gamma=0.2$ reaches 19.04\,dB, against
17.82\,dB ungated. In a run whose clean pool also contained noisier Vaani speech, the gate instead lowered
SI-SDR (10.5 vs.\ 16.5\,dB ungated), consistent with activity targets that mislabel real background noise as
speech.
\item \textbf{Noise level, not duration, limits performance.} On the scored test clips, per-clip SI-SDR correlates
0.82 with the estimated input SNR and 0.09 with clip length. Clips below $-6$\,dB input SNR (39\%) average
8.5\,dB, and clips above 0\,dB average 20.7\,dB. On the pseudo-test our ensemble (13.36\,dB) already exceeds the
ideal ratio mask (12.99\,dB) and approaches the ideal phase-sensitive mask (15.15\,dB): mask-style enhancement has
little headroom left at such SNRs.
\item \textbf{Pseudo-test fidelity.} The estimator predicted 14.52 for our final entry, which scored 14.65.
Validation-holdout estimates were optimistic by about 5\,dB, because test clips are 2.7$\times$ longer and much
noisier.
\end{itemize}

\textbf{Novelty.} Our contributions are (i) an event-gated BSRNN that learns event activity jointly with
separation; (ii) ASR-consistency training through the frozen mandated ASR; (iii) wall-clock-scheduled training with
validation-matched on-the-fly mixing; and (iv) a holdout-derived pseudo-test together with a calibrated
estimate of the leaderboard score, used to choose submissions.

\textbf{Conclusion.} A single learned, event-aware enhancer trained on statistically matched mixtures delivers
strong SI-SDR, with a positive ASR gain on both natural and synthetic speech. The remaining gap lies in very low-SNR
clips, where stronger separation models and longer test-condition training are the clearest next steps.

\bibliographystyle{IEEEtran}
\begin{thebibliography}{10}\setlength{\itemsep}{0pt}\footnotesize
\bibitem{vaani} S.~Pulikodan \emph{et~al.}, ``VAANI: Capturing the language landscape for an inclusive digital India,'' arXiv:2603.28714, 2026.
\bibitem{demucs} A.~D\'efossez, G.~Synnaeve, and Y.~Adi, ``Real time speech enhancement in the waveform domain,'' in \emph{Proc. Interspeech}, 2020.
\bibitem{bsrnn} Y.~Luo and J.~Yu, ``Music source separation with band-split RNN,'' \emph{IEEE/ACM TASLP}, vol.~31, 2023.
\bibitem{sisdr} J.~Le~Roux, S.~Wisdom, H.~Erdogan, and J.~R.~Hershey, ``SDR -- half-baked or well done?'' in \emph{Proc. ICASSP}, 2019.
\bibitem{pwg} R.~Yamamoto, E.~Song, and J.-M.~Kim, ``Parallel WaveGAN,'' in \emph{Proc. ICASSP}, 2020.
\bibitem{esc50} K.~J.~Piczak, ``ESC: Dataset for environmental sound classification,'' in \emph{Proc. ACM Multimedia}, 2015.
\bibitem{fleurs} A.~Conneau \emph{et~al.}, ``FLEURS: Few-shot learning evaluation of universal representations of speech,'' in \emph{Proc. IEEE SLT}, 2022.
\bibitem{sravaani} ARTPARK-IISc, ``SraVaani-1.0 ASR,'' \url{https://huggingface.co/ARTPARK-IISc/SraVaani-1.0}, 2026.
\bibitem{tdt} H.~Xu \emph{et~al.}, ``Efficient sequence transduction by jointly predicting tokens and durations,'' in \emph{Proc. ICML}, 2023.
\bibitem{adamw} I.~Loshchilov and F.~Hutter, ``Decoupled weight decay regularization,'' in \emph{Proc. ICLR}, 2019.
\end{thebibliography}

\end{document}
